% Generated from locale/tr/content/turing-machines/undecidability/unsolvability-decision-problem.tex; do not hand-edit.


\olfileid[tr]{tur}{und}{uns} 
\olsection{Karar Problemi Çözülemezdir}

\begin{thm}
\ollabel{thm:decision-prob}
Karar problemi çözülemezdir: Öyle bir Turing makinesi~$D$ yoktur ki birinci
dereceden mantığın !!a{sentence}~$!B$ şeridinde girdi olarak verildiğinde
$D$ sonunda dursun ve ancak ve ancak $!B$ geçerliyse $1$, aksi durumda $0$
çıktısını versin.
\end{thm}

\begin{proof}
Karar probleminin çözülebilir, yani bir $D$ Turing makinesinin var olduğunu
varsayalım. O zaman durma problemini şöyle çözebilirdik. Girdi olarak $M_e$
Turing makinesinin $e$ numarasıyla $w$ girdisi verildiğinde karşılık gelen
!!{sentence}~$!T(M_e,w)\lif !E(M_e,w)$'yi hesaplayıp şeritteki en soldaki
kareyi tararken duran bir $E$ Turing makinesi kurarız. $E\concat D$ makinesi,
$e$ ile $w$ girdisi verildiğinde önce $!T(M_e,w)\lif !E(M_e,w)$'yi hesaplar,
sonra karar problemi makinesi $D$'yi bu girdide çalıştırırdı. $D$ ancak ve
ancak $!T(M_e,w)\lif !E(M_e,w)$ geçerliyse $1$, aksi durumda $0$ çıktısıyla
durur. \olref[ver]{lem:halt-if-valid} ile
\olref[ver]{lem:valid-if-halt} gereği $!T(M_e,w)\lif !E(M_e,w)$ ancak ve
ancak $M_e$, $w$ girdisinde durursa geçerlidir. Böylece $E\concat D$, $e$ ve
$w$ girdisi verildiğinde ancak ve ancak $M_e$, $w$ girdisinde durursa $1$,
aksi durumda $0$ çıktısıyla dururdu. Başka bir deyişle $E\concat D$ durma
problemini çözerdi. Ancak \olref[hal]{thm:halting-problem} gereği böyle bir
Turing makinesinin bulunamayacağını biliyoruz.
\end{proof}

\begin{cor}\ollabel{cor:undecidable-sat}%
Birinci dereceden mantığın keyfî bir !!{sentence}{}nin doyurulabilir olup
olmadığı karar verilemezdir.
\end{cor}

\begin{proof}
  Doyurulabilirliğe bir $S$ Turing makinesince karar verilebildiğini
  varsayalım. O zaman karar problemini şöyle çözebilirdik: Girdi olarak
  !!a{sentence}~$B$ verildiğinde $!B$'yi bir kare sağa taşıyın. $1$. kareye
  dönüp $\lnot$ simgesini yazın.

  Şimdi $S$ Turing makinesini çalıştırın. Makine sonunda şeritte ya $1$
  ($\lnot !B$ doyurulabilir ise) ya da $0$ ($\lnot !B$ doyurulamaz ise)
  çıktısıyla durur. $1$. karede bir~$\TMstroke$ varsa onu silin; $1$. kare
  boşsa bir~$\TMstroke$ yazıp durun.

  Bu Turing makinesi her zaman durur ve ancak ve ancak $\lnot !B$
  doyurulamazsa $1$, aksi durumda $0$ çıktısını verir. $!B$ ancak ve ancak
  $\lnot !B$ doyurulamazsa geçerli olduğundan makine ancak ve ancak $!B$
  geçerliyse $1$, aksi durumda $0$ çıktısını verir; yani karar problemini
  çözerdi.
\end{proof}

\begin{explain}
Dolayısıyla ``$!B$, birinci dereceden mantığın geçerli bir !!{sentence} midir?''
sorusuna her zaman doğru bir ``evet'' ya da ``hayır'' yanıtı veren hiçbir
Turing makinesi yoktur. Bununla birlikte her zaman doğru bir ``evet'' yanıtı
veren, ama yanıt ``hayır'' ise yalnızca durmayan bir Turing makinesi
\emph{vardır}. Bu, birinci dereceden mantığın sağlamlık ve tamlık teoremleriyle
!!{derivation}s etkili biçimde numaralandırılabilmesi olgularından çıkar.
\end{explain}

\begin{thm}
  \ollabel{thm:valid-ce}%
  Birinci dereceden !!{sentence}s geçerliliği yarı karar verilebilirdir:
  Bir $E$ Turing makinesi vardır: Birinci dereceden mantığın
  !!a{sentence}~$!B$ şeridinde girdi olarak
  verildiğinde $E$ ancak ve ancak $!B$ geçerliyse sonunda durup $1$ çıktısını
  verir; aksi durumda durmaz.
\end{thm}

\begin{proof}
  Birinci dereceden mantığın bütün olası !!{derivation}s etkili bir
  algoritmayla birbiri ardına üretilebilir. $E$ makinesi bunu yapar ve
  $\Proves !B$ olduğunu gösteren !!a{derivation} bulduğunda $1$ çıktısıyla
  durur. Sağlamlık teoremine göre $E$, $1$ çıktısıyla durursa bunun nedeni
  $\Entails !B$ olmasıdır. Tamlık teoremine göre $\Entails !B$ ise
  $\Proves !B$ olduğunu gösteren !!a{derivation} vardır. $E$ bütün olası
  !!{derivation}s sistematik olarak ürettiğinden sonunda $\Proves !B$
  olduğunu gösteren birini bulacak ve $1$ çıktısıyla duracaktır.
\end{proof}

